% \documentclass[polish, 8pt]{article}
\documentclass[polish, 10pt]{extarticle}
\usepackage[top=2cm, bottom=2cm, left=3cm, right=3cm]{geometry}
\usepackage{graphicx} % Required for inserting images

\usepackage[polish]{babel}
\usepackage{polski}
\usepackage{placeins}
\usepackage{dirtytalk}
% Set page size and margins
% Replace `letterpaper' with `a4paper' for UK/EU standard size
\usepackage{csquotes}
% Useful packages
\usepackage{float}
\usepackage{amsmath}
\usepackage{graphicx}
\usepackage{hyperref}
\usepackage{graphicx}
\usepackage{listings}
\usepackage{xcolor}
\usepackage{amssymb}
\usepackage{enumitem}
\usepackage{float}
\usepackage{amsmath}
\usepackage{amssymb}
\usepackage{amsthm}

\definecolor{codegreen}{rgb}{0,0.6,0}
\definecolor{codegray}{rgb}{0.5,0.5,0.5}
\definecolor{codepurple}{rgb}{0.58,0,0.82}
\definecolor{backcolour}{rgb}{0.95,0.95,0.92}

\newtheorem*{twierdzenie}{Twierdzenie}
\theoremstyle{definition}
\newtheorem*{definicja}{Definicja}

\lstdefinestyle{mystyle}{
    backgroundcolor=\color{backcolour},   
    commentstyle=\color{codegreen},
    keywordstyle=\color{magenta},
    numberstyle=\tiny\color{codegray},
    stringstyle=\color{codepurple},
    basicstyle=\ttfamily\footnotesize,
    breakatwhitespace=false,         
    breaklines=true,                 
    captionpos=b,                    
    keepspaces=true,                 
    numbers=left,                    
    numbersep=5pt,                  
    showspaces=false,                
    showstringspaces=false,
    showtabs=false,                  
    tabsize=2
}
\lstset{style=mystyle}
\hypersetup{
    colorlinks=true,
    linktoc=all     %set to all if you want both sections and subsections linked
}
\usepackage{mdframed}
\usepackage{titling}
\setlength{\droptitle}{-2cm} 

\begin{document}

\title{Kolokwium I -- Zestaw A}
\author{}
\date{}

\maketitle
\vspace{-2cm}

\section*{Zadanie 1 (25 pkt)}
% W kapeluszu znajdują się dwie monety. Monety wyglądają tak samo, ale jedna jest uczciwa (z prawdopodobieństwem wypadnięcia orła równym \(1/2\)), natomiast druga jest stronnicza, z prawdopodobieństwem orła równym \(1/4\). Jedną z monet losowo wyciąga się z kapelusza, nie wiedząc, która to. Wybraną monetę nazywamy \emph{monetą \(C\)}.

% \begin{enumerate}
%   \item Monetę \(C\) rzucono dwukrotnie i za każdym razem wypadł orzeł. Znając tę informację, jakie jest prawdopodobieństwo, że moneta \(C\) jest uczciwa?
%   \item Czy zdarzenia „pierwszy rzut monety \(C\) daje orła” oraz „drugi rzut monety \(C\) daje orła” są niezależne? Uzasadnij.
%   \item Wyznacz prawdopodobieństwo, że w \(10\) rzutach monety \(C\) pojawią się dokładnie \(3\) orły. (Moneta jest jednakowo prawdopodobnie jedną z dwóch monet; nie zakładaj, że wcześniej wypadły dwa orły jak w punkcie (a).)
% \end{enumerate}
W kapeluszu znajdują się dwie monety, których nie można odróżnić wzrokiem. Jedna z nich jest uczciwa (prawdopodobieństwo orła $1/2$), a druga obciążona (prawdopodobieństwo orła $1/4$). Losujemy z kapelusza jedną monetę – nazwijmy ją monetą $C$.
\begin{enumerate}
\item Rzucono dwukrotnie monetą $C$ i za każdym razem wypadł orzeł. Jakie jest prawdopodobieństwo, że wylosowana moneta jest uczciwa?
\item Czy zdarzenia: „w pierwszym rzucie monetą $C$ wypadł orzeł” oraz „w drugim rzucie monetą $C$ wypadł orzeł” są niezależne? Uzasadnij odpowiedź.
\item Wracamy do momentu wyboru monety (zapominając o wynikach z punktu a). Jakie jest prawdopodobieństwo uzyskania dokładnie $3$ orłów w $5$ rzutach wylosowaną monetą $C$?
\end{enumerate}
\section*{Rozwiązanie}

\textbf{Oznaczenia:}
\begin{itemize}
    \item $M_U$ – wylosowano monetę uczciwą, $P(M_U) = \frac{1}{2}$, $P(O|M_U) = \frac{1}{2}$.
    \item $M_O$ – wylosowano monetę obciążoną, $P(M_O) = \frac{1}{2}$, $P(O|M_O) = \frac{1}{4}$.
\end{itemize}

\begin{enumerate}[label=\alph*)]
    \item \textbf{Prawdopodobieństwo, że moneta jest uczciwa (Tw. Bayesa)}
    
    Niech $2O$ oznacza wyrzucenie dwóch orłów. Obliczamy $P(M_U|2O)$.
    
    Z prawdopodobieństwa całkowitego:
    \[
    P(2O) = P(2O|M_U)P(M_U) + P(2O|M_O)P(M_O)
    \]
    \[
    P(2O) = \left(\frac{1}{2}\right)^2 \cdot \frac{1}{2} + \left(\frac{1}{4}\right)^2 \cdot \frac{1}{2} = \frac{1}{8} + \frac{1}{32} = \frac{5}{32}
    \]
    
    Z tw. Bayesa:
    \[
    P(M_U|2O) = \frac{P(2O|M_U)P(M_U)}{P(2O)} = \frac{\frac{1}{8}}{\frac{5}{32}} = \frac{1}{8} \cdot \frac{32}{5} = \mathbf{\frac{4}{5}}
    \]

    \item \textbf{Niezależność zdarzeń w dwóch rzutach}
    
    Niech $A$ – orzeł w 1. rzucie, $B$ – orzeł w 2. rzucie.
    \begin{itemize}
        \item $P(A \cap B) = P(2O) = \frac{5}{32}$ (z punktu a).
        \item $P(A) = P(O|M_U)P(M_U) + P(O|M_O)P(M_O) = \frac{1}{2}\cdot\frac{1}{2} + \frac{1}{4}\cdot\frac{1}{2} = \frac{3}{8}$.
        \item Z symetrii $P(B) = P(A) = \frac{3}{8}$.
    \end{itemize}
    Sprawdzamy warunek niezależności $P(A \cap B) \stackrel{?}{=} P(A) \cdot P(B)$:
    \[
    \frac{5}{32} \neq \left(\frac{3}{8}\right)^2 \quad \Rightarrow \quad \frac{10}{64} \neq \frac{9}{64}
    \]
    \textbf{Odpowiedź:} Zdarzenia \textbf{nie są niezależne}.

    \item \textbf{3 orły w 5 rzutach (Schemat Bernoulliego)}
    
    Niech $S_5$ oznacza liczbę orłów w 5 rzutach. Szukamy $P(S_5=3)$.
    Korzystamy ze wzoru na prawdopodobieństwo całkowite, gdzie $P_n(k) = \binom{n}{k}p^k(1-p)^{n-k}$:
    
    \[
    P(S_5=3) = P(S_5=3|M_U)P(M_U) + P(S_5=3|M_O)P(M_O)
    \]
    
    \[
    P(S_5=3|M_U) = \binom{5}{3}\left(\frac{1}{2}\right)^5 = 10 \cdot \frac{1}{32} = \frac{5}{16}
    \]
    
    \[
    P(S_5=3|M_O) = \binom{5}{3}\left(\frac{1}{4}\right)^3\left(\frac{3}{4}\right)^2 = 10 \cdot \frac{1}{64} \cdot \frac{9}{16} = \frac{45}{512}
    \]
    
    Podstawiamy:
    \[
    P(S_5=3) = \frac{5}{16} \cdot \frac{1}{2} + \frac{45}{512} \cdot \frac{1}{2} = \frac{5}{32} + \frac{45}{1024} = \frac{160 + 45}{1024} = \mathbf{\frac{205}{1024}}
    \]
\end{enumerate}

\section*{Zadanie 2 (15 pkt)}
Niech \(A\) i \(B\) będą zdarzeniami losowymi.
Różnicą symetryczną \(A \triangle B\) nazywamy zdarzenie zachodzące wtedy, gdy zachodzi \(A\) lub zachodzi \(B\), ale nie oba jednocześnie (w logice odpowiada to alternatywie wykluczającej \emph{XOR}). Korzystając z definicji i własności prawdopodobieństwa, wykaż, że:
\begin{enumerate}
\item $P(A \triangle B) \;=\; P(A) + P(B) - 2P(A \cap B)$
\item $P(A) - P(B) \leq P(A \triangle B)$
\end{enumerate}
\section*{Rozwiązanie}

\subsection*{Dowód punktu 1.}
Z definicji różnicy symetrycznej wiemy, że:
\[ A \triangle B = (A \setminus B) \cup (B \setminus A) \]
Zbiory te są rozłączne, tzn. $(A \setminus B) \cap (B \setminus A) = \emptyset$, więc:
\begin{equation}
    P(A \triangle B) = P(A \setminus B) + P(B \setminus A) \label{eq:sym_diff_sum}
\end{equation}
Rozważmy teraz prawdopodobieństwo różnicy zbiorów. Wiemy, że $A = (A \setminus B) \cup (A \cap B)$, gdzie składniki sumy są rozłączne. Stąd:
\[ P(A) = P(A \setminus B) + P(A \cap B) \]
Przekształcając powyższe równanie, otrzymujemy wzory na prawdopodobieństwa różnic:
\begin{align*}
    P(A \setminus B) &= P(A) - P(A \cap B) \\
    P(B \setminus A) &= P(B) - P(A \cap B)
\end{align*}
Podstawiając te zależności do równania (\ref{eq:sym_diff_sum}), otrzymujemy tezę:
\begin{align*}
    P(A \triangle B) &= \bigl( P(A) - P(A \cap B) \bigr) + \bigl( P(B) - P(A \cap B) \bigr) \\
    &= P(A) + P(B) - 2P(A \cap B)
\end{align*}
\hfill $\square$

\subsection*{Dowód punktu 2.}
Zauważmy, że iloczyn zdarzeń jest podzbiorem każdego z nich, tzn. $(A \cap B) \subseteq B$. Z własności prawdopodobieństwa wynika, że:
\[ P(A \cap B) \leq P(B) \]
Mnożąc tę nierówność obustronnie przez $-2$, otrzymujemy:
\[ -2P(A \cap B) \geq -2P(B) \]
Dodajmy teraz do obu stron wyrażenie $P(A) + P(B)$:
\[ P(A) + P(B) - 2P(A \cap B) \geq P(A) + P(B) - 2P(B) \]
Lewa strona to dokładnie $P(A \triangle B)$ (na podstawie punktu 1.), natomiast prawa może być uproszczona:
\[ P(A \triangle B) \geq P(A) - P(B) \]
Co kończy dowód.
\hfill $\square$

\section*{Zadanie 3 (10 pkt)}
Podaj definicję rozkładu prawdopodobieństwa zmiennej losowej. Sformułuj i udowodnij twierdzenie mówiące, że ten rozkład jest miarą probabilistyczną.

\subsection*{Rozwiązanie}
Slajdy 2.2.1, 2.2.3 i 2.2.4 z wykładu.

\section*{Zadanie 4 (30 pkt)}

Niech $X$ będzie zmienną losową o nieznanym parametrze $\theta \in \mathbb{R}_+$, której funkcja
gęstości prawdopodobieństwa ma postać
\[
f_\theta(x) =
\begin{cases}
\dfrac{2x}{\theta^2}, & \text{dla } 0 \leq x \leq \theta,\\[4pt]
0, & \text{w przeciwnym przypadku.}
\end{cases}
\]

\begin{enumerate}[label=\alph*)]
    \item Dla jakich wartości parametru $\theta$ funkcja $f_\theta$ jest
    funkcją gęstości prawdopodobieństwa zmiennej losowej $X$?

    \item Wyznacz dystrybuantę $F_X(x)$ zmiennej losowej $X$.

    \item Oblicz wartość oczekiwaną $E[2X+3]$, $Var(X)$ oraz medianę
    zmiennej losowej $X$ w zależności od parametru $\theta$.

    \item Oblicz prawdopodobieństwo $P\bigl(X^2 > 1\bigr)$ w zależności od parametru $\theta$.
\end{enumerate}

\section*{Rozwiązanie}

\subsection*{a) Wyznaczenie wartości parametru $\theta$}
Aby funkcja $f_\theta(x)$ była funkcją gęstości, musi spełniać dwa warunki:
\begin{enumerate}
    \item $f_\theta(x) \geq 0$ dla każdego $x \in \mathbb{R}$.
    \item $\int_{-\infty}^{\infty} f_\theta(x) \, dx = 1$.
\end{enumerate}

\noindent \textbf{Warunek 1:} Dla $\theta > 0$ oraz $x \in [0, \theta]$, wyrażenie $\frac{2x}{\theta^2}$ jest nieujemne. \\
\textbf{Warunek 2:} Obliczamy całkę:
\[
\int_{-\infty}^{\infty} f_\theta(x) \, dx = \int_0^\theta \frac{2x}{\theta^2} \, dx = \frac{2}{\theta^2} \int_0^\theta x \, dx
\]
\[
= \frac{2}{\theta^2} \cdot \left[ \frac{x^2}{2} \right]_0^\theta = \frac{2}{\theta^2} \cdot \left( \frac{\theta^2}{2} - 0 \right) = 1
\]
Całka wynosi 1 dla każdego $\theta \neq 0$. Z treści zadania wiemy, że $\theta \in \mathbb{R}_+$.

\noindent \textbf{Odpowiedź:} Funkcja jest gęstością dla każdego $\theta > 0$.

\subsection*{b) Wyznaczenie dystrybuanty $F_X(x)$}
Dystrybuantę wyznaczamy ze wzoru $F_X(x) = \int_{-\infty}^x f_\theta(t) \, dt$.

\begin{itemize}
    \item Dla $x < 0$: $F_X(x) = 0$.
    \item Dla $0 \leq x \leq \theta$:
    \[
    F_X(x) = \int_0^x \frac{2t}{\theta^2} \, dt = \frac{2}{\theta^2} \cdot \frac{x^2}{2} = \frac{x^2}{\theta^2}
    \]
    \item Dla $x > \theta$: $F_X(x) = 1$.
\end{itemize}

\noindent \textbf{Odpowiedź:}
\[
F_X(x) =
\begin{cases}
0, & \text{dla } x < 0, \\[6pt]
\dfrac{x^2}{\theta^2}, & \text{dla } 0 \leq x \leq \theta, \\[6pt]
1, & \text{dla } x > \theta.
\end{cases}
\]

\subsection*{c) Wartość oczekiwana, wariancja i mediana}

\textbf{1. Wartość oczekiwana $E[X]$:}
\[
E[X] = \int_0^\theta x \cdot \frac{2x}{\theta^2} \, dx = \frac{2}{\theta^2} \int_0^\theta x^2 \, dx = \frac{2}{\theta^2} \left[ \frac{x^3}{3} \right]_0^\theta = \frac{2}{\theta^2} \cdot \frac{\theta^3}{3} = \frac{2}{3}\theta
\]
Stąd wartość oczekiwana $E[2X+3]$:
\[
E[2X+3] = 2E[X] + 3 = 2\left(\frac{2}{3}\theta\right) + 3 = \frac{4}{3}\theta + 3
\]

\textbf{2. Wariancja $Var(X)$:}
Najpierw obliczamy drugi moment zwykły $E[X^2]$:
\[
E[X^2] = \int_0^\theta x^2 \cdot \frac{2x}{\theta^2} \, dx = \frac{2}{\theta^2} \int_0^\theta x^3 \, dx = \frac{2}{\theta^2} \left[ \frac{x^4}{4} \right]_0^\theta = \frac{\theta^2}{2}
\]
Obliczamy wariancję ze wzoru $Var(X) = E[X^2] - (E[X])^2$:
\[
Var(X) = \frac{\theta^2}{2} - \left(\frac{2}{3}\theta\right)^2 = \frac{\theta^2}{2} - \frac{4\theta^2}{9} = \frac{9\theta^2 - 8\theta^2}{18} = \frac{\theta^2}{18}
\]

\textbf{3. Mediana $Me$:}
Szukamy takiej wartości $Me$, dla której $F_X(Me) = 0,5$:
\[
\frac{Me^2}{\theta^2} = \frac{1}{2} \implies Me^2 = \frac{\theta^2}{2} \implies Me = \frac{\theta}{\sqrt{2}} = \frac{\sqrt{2}}{2}\theta
\]

\subsection*{d) Prawdopodobieństwo $P(X^2 > 1)$}
Nierówność $X^2 > 1$ jest równoważna $X > 1$ (ponieważ $X \ge 0$).
\[
P(X > 1) = 1 - P(X \leq 1) = 1 - F_X(1)
\]
Rozwiązanie zależy od wartości parametru $\theta$:
\begin{itemize}
    \item Jeśli $\theta \leq 1$, to $X$ przyjmuje wartości tylko z przedziału $[0, \theta]$, więc $P(X > 1) = 0$.
    \item Jeśli $\theta > 1$, to $1$ należy do nośnika funkcji, więc:
    \[
    P(X > 1) = 1 - \frac{1^2}{\theta^2} = 1 - \frac{1}{\theta^2}
    \]
\end{itemize}

\noindent \textbf{Odpowiedź:}
\[
P(X^2 > 1) =
\begin{cases}
0, & \text{gdy } \theta \leq 1, \\[6pt]
1 - \dfrac{1}{\theta^2}, & \text{gdy } \theta > 1.
\end{cases}
\]

\section*{Zadanie 5 (20 pkt)}
Zmienna losowa $X$ ma rozkład jednostajny $U(-2, 3)$.
Wyznacz rozkład zmiennej losowej $Y = (X-1)^2$.

\section*{Rozwiązanie}

\subsection*{1. Rozkład zmiennej losowej \(X\)}
Zmienna losowa \(X\) ma rozkład jednostajny na przedziale \((-2, 3)\). Jej gęstość prawdopodobieństwa wyraża się wzorem:
\begin{equation}
    f_X(x) = \begin{cases} 
        \frac{1}{3 - (-2)} = \frac{1}{5} & \text{dla } x \in (-2, 3) \\
        0 & \text{dla } x \notin (-2, 3)
    \end{cases}
\end{equation}

\subsection*{2. Wyznaczenie zbioru wartości zmiennej \(Y\)}
Funkcja transformująca to \(g(x) = (x-1)^2\).
\begin{itemize}
    \item Wierzchołek paraboli znajduje się w punkcie \(x=1\), który należy do przedziału \((-2, 3)\), zatem minimalna wartość to \(Y_{\min} = 0\).
    \item Sprawdzamy krańce przedziału:
    \begin{itemize}
        \item dla \(x = 3\): \(y = (3-1)^2 = 4\)
        \item dla \(x = -2\): \(y = (-2-1)^2 = 9\)
    \end{itemize}
\end{itemize}
Zatem zmienna \(Y\) przyjmuje wartości z przedziału \([0, 9)\).

\subsection*{3. Wyznaczenie gęstości \(f_Y(y)\) metodą dystrybuanty}
Z definicji dystrybuanty:
\[ F_Y(y) = P((X-1)^2 \le y) = P(1-\sqrt{y} \le X \le 1+\sqrt{y}) \]
\[ F_Y(y) = F_X(1+\sqrt{y}) - F_X(1-\sqrt{y}) \]
Różniczkujemy powyższe równanie obustronnie względem \(y\), aby otrzymać gęstość \(f_Y(y)\). Korzystamy z pochodnej funkcji złożonej oraz faktu, że \(F_X'(x) = f_X(x)\):

\begin{align*}
    f_Y(y) = \frac{d}{dy} F_Y(y) &= \frac{d}{dy} \left( F_X(1+\sqrt{y}) - F_X(1-\sqrt{y}) \right) \\[1em]
    &= f_X(1+\sqrt{y}) \cdot (1+\sqrt{y})' - f_X(1-\sqrt{y}) \cdot (1-\sqrt{y})' \\[1em]
    &= f_X(1+\sqrt{y}) \cdot \frac{1}{2\sqrt{y}} - f_X(1-\sqrt{y}) \cdot \left( -\frac{1}{2\sqrt{y}} \right) \\[1em]
    &= \frac{1}{2\sqrt{y}} \left[ f_X(1+\sqrt{y}) + f_X(1-\sqrt{y}) \right]
\end{align*}

\subsection*{4. Analiza przedziałów dla \(Y\)}
Musimy sprawdzić, kiedy argumenty $f_X$ należą do przedziału \((-2, 3)\). Do rozważenia mamy 2 argumenty: $(1+\sqrt{y})$ i $(1-\sqrt{y})$:. Pierwszy przypadek $(1+\sqrt{y})$:

\begin{equation*}
-2 < 1 + \sqrt{y} < 3
\end{equation*}
\begin{equation*}
-3 < \sqrt{y} < 2
\end{equation*}
\begin{equation*}
0 < y < 4
\end{equation*}
Drugi przypadek $(1-\sqrt{y})$:
\begin{equation*}
-2 < 1 - \sqrt{y} < 3
\end{equation*}
\begin{equation*}
-3 < - \sqrt{y} < 2
\end{equation*}
\begin{equation*}
-2 < \sqrt{y} < 3
\end{equation*}
\begin{equation*}
0 < y < 9
\end{equation*}
Z ogólnego wzoru
\[
    f_Y(y) = \frac{1}{2\sqrt{y}} \left[ f_X(1+\sqrt{y}) + f_X(1-\sqrt{y}) \right]
\]
korzystamy teraz z faktu, że
\[
    f_X(x) = \begin{cases}
        \frac{1}{5}, & x \in (-2,3),\\
        0, & x \notin (-2,3).
    \end{cases}
\]
Z analizy przedziałów mamy:
\begin{itemize}
    \item dla \(0 < y < 4\) zachodzi jednocześnie
    \[
        1+\sqrt{y} \in (-2,3) \quad \text{oraz} \quad 1-\sqrt{y} \in (-2,3),
    \]
    więc
    \[
        f_X(1+\sqrt{y}) = \frac{1}{5}, \qquad f_X(1-\sqrt{y}) = \frac{1}{5},
    \]
    a zatem
    \[
        f_Y(y) = \frac{1}{2\sqrt{y}}\left( \frac{1}{5} + \frac{1}{5} \right)
        = \frac{1}{2\sqrt{y}} \cdot \frac{2}{5}
        = \frac{1}{5\sqrt{y}}.
    \]
    \item dla \(4 \le y < 9\) mamy
    \[
        1+\sqrt{y} \notin (-2,3) \quad (\text{bo } 1+\sqrt{y} \ge 3),
    \]
    więc
    \[
        f_X(1+\sqrt{y}) = 0,
    \]
    natomiast wciąż
    \[
        1-\sqrt{y} \in (-2,3),
    \]
    czyli
    \[
        f_X(1-\sqrt{y}) = \frac{1}{5}.
    \]
    Wtedy
    \[
        f_Y(y) = \frac{1}{2\sqrt{y}}\left( 0 + \frac{1}{5} \right)
        = \frac{1}{10\sqrt{y}}.
    \]
    \item ponadto z części 2 wiemy, że zmienna \(Y\) przyjmuje wartości w przedziale \([0,9]\), więc na zewnątrz tego przedziału gęstość również jest równa zero.
\end{itemize}

\section*{Odpowiedź}
Gęstość rozkładu zmiennej losowej $Y$:
\begin{equation*}
    f_Y(y) = \begin{cases} 
        \frac{1}{5\sqrt{y}} & \text{dla } y \in (0, 4) \\[1em]
        \frac{1}{10\sqrt{y}} & \text{dla } y \in [4, 9) \\[1em]
        0 & \text{w p.p.}
    \end{cases}
\end{equation*}


\end{document}

